% Using ACM conference document class for professional formatting
\documentclass[sigconf,nonacm]{acmart}

% Including essential packages for algorithms and formatting
\usepackage{algorithm}
\usepackage{algpseudocode}
\usepackage{amsmath}
\usepackage{booktabs}
\usepackage{xcolor}

% Defining custom commands for clarity
\newcommand{\openshmem}{\texttt{OpenSHMEM}}
\newcommand{\mlir}{\texttt{MLIR}}
\newcommand{\barrier}{\texttt{openshmem.barrier}}

% Setting up metadata for the document
\AtBeginDocument{%
  \providecommand\BibTeX{{%
    Bib\TeX}}}

\setcopyright{none}
\copyrightyear{2025}
\acmYear{2025}
\acmDOI{}

% Starting the document
\begin{document}

% Providing a title for context
\title{Sync Removal Optimization for OpenSHMEM in MLIR}

% Omitting author for simplicity; can be added as needed
\author{}
\date{}

% Generating the algorithm figure
\begin{algorithm}
\caption{Synchronization Removal Pass for \openshmem{} Dialect in \mlir}
\label{alg:sync-removal}
\begin{algorithmic}[1]
\Require{\openshmem{} IR containing synchronization operations}
\Ensure{Optimized IR with redundant synchronizations removed}
\Statex % Spacer for clarity
\State $ops \gets$ Collect all \barrier{} operations in the IR \Comment{Identify synchronization points}
\For{each $op \in ops$} \Comment{Iterate over barriers}
    \State $deps \gets$ Compute memory dependencies via \texttt{MemoryDependenceAnalysis} \Comment{Check data conflicts}
    \State $control \gets$ Analyze control flow dependencies for $op$ \Comment{Verify control constraints}
    \If{$deps$ contains no read-write or write-write conflicts across $op$} \Comment{No data hazards}
        \If{$control$ imposes no critical dependencies on $op$} \Comment{No control hazards}
            \State Remove $op$ using \texttt{op.erase()} \Comment{Eliminate redundant barrier}
        \EndIf
    \EndIf
\EndFor
\State Simplify IR using \texttt{Canonicalizer} \Comment{Clean up residual operations}
\State \Return Optimized IR \Comment{Output optimized program}
\end{algorithmic}
\end{algorithm}

% Ending the document
\end{document}
